\section{Absorption models and the method of characteristics}
\label{sec:characteristics}

We now consider a population split between the exterior and interior of a single
compartment.  Transfer moves an exterior particle to the interior; interior
particles may then die or reproduce.  The models are deliberately elementary, but
their two-dimensional state and linear per-capita rates already produce the
generating-function PDEs needed later for systems with rupture.

Let \(X_t\) and \(Y_t\) denote the interior and exterior counts, respectively, and
write
\begin{equation}
  p_{i,j}(t)=\Prob(X_t=i,Y_t=j),
  \qquad
  G(x,y,t)=\E[x^{X_t}y^{Y_t}]
  =\sum_{i,j\geq0}p_{i,j}(t)x^iy^j.
  \label{eq:absorption-pgf-definition}
\end{equation}
Every model starts from \(X_0=0\), \(Y_0=N\), and hence
\begin{equation}
  G(x,y,0)=y^N.
  \label{eq:absorption-initial-pgf}
\end{equation}
Three cases are treated in increasing order of difficulty.  Transfer alone and
transfer followed by death have finite reachable state spaces and admit direct
single-particle solutions.  Adding birth opens an infinite state space.  The
branching property still raises the one-founder generating function to the power
\(N\), but it no longer supplies that one-founder function; the PDE must then be
solved.

\subsection{Transfer only}
\label{sec:transfer-only}

Each exterior particle transfers at rate \(\alpha>0\), so the only jump is
\begin{equation}
  (i,j)\longrightarrow(i+1,j-1)
  \quad\text{at rate }\alpha j.
  \label{eq:transfer-only-rate}
\end{equation}
Equivalently, over a short interval \(h\),
\begin{align}
  \Prob(\Delta X=1,\Delta Y=-1\mid X_t=i,Y_t=j)
    &=\alpha jh+o(h),
  \notag\\
  \Prob(\Delta X=0,\Delta Y=0\mid X_t=i,Y_t=j)
    &=1-\alpha jh+o(h).
  \label{eq:transfer-only-short-time}
\end{align}
The componentwise forward equation is
\begin{equation}
  \frac{\dd p_{i,j}}{\dd t}
  =\alpha(j+1)p_{i-1,j+1}-\alpha jp_{i,j}.
  \label{eq:transfer-only-forward}
\end{equation}

The generating-function PDE follows directly from the generator.  Applied to the
monomial \(f(i,j)=x^iy^j\),
\begin{align}
  \mathcal Lf(i,j)
  &=\alpha j\{x^{i+1}y^{j-1}-x^iy^j\}
  \notag\\
  &=\alpha(x-y)\frac{\partial}{\partial y}(x^iy^j).
\end{align}
Taking expectations gives
\begin{equation}
  \frac{\partial G}{\partial t}
  =\alpha(x-y)\frac{\partial G}{\partial y},
  \qquad G(x,y,0)=y^N.
  \label{eq:transfer-only-pde}
\end{equation}
This generator calculation is equivalent to multiplying
\cref{eq:transfer-only-forward} by \(x^iy^j\), summing and shifting indices, but it
keeps the jump that produces each derivative visible.

The direct solution uses independence.  For one initial exterior particle the only
states are \((0,1)\) and \((1,0)\), with
\begin{equation}
  p_{0,1}(t)=\e^{-\alpha t},
  \qquad
  p_{1,0}(t)=1-\e^{-\alpha t}.
  \label{eq:transfer-only-states}
\end{equation}
Thus
\begin{equation}
  G_1(x,y,t)=(1-\e^{-\alpha t})x+\e^{-\alpha t}y.
\end{equation}
The \(N\) initial particles move independently, so
\begin{equation}
  \boxed{
  G_N(x,y,t)=
  \left\{(1-\e^{-\alpha t})x+\e^{-\alpha t}y\right\}^{N}.}
  \label{eq:G-transfer-only}
\end{equation}
In particular,
\begin{equation}
  X_t\sim\operatorname{Binomial}(N,1-\e^{-\alpha t}),
  \qquad Y_t=N-X_t.
  \label{eq:transfer-binomial}
\end{equation}
Direct differentiation confirms that \cref{eq:G-transfer-only} satisfies
\cref{eq:transfer-only-pde}.

\begin{figure}[htbp]
  \centering
  \includegraphics[width=0.58\textwidth]{abs1.pdf}
  \caption{State probabilities for one particle in the transfer-only model.  The
  probability of remaining outside decays at rate \(\alpha\); the complementary
  probability records transfer to the interior.}
  \label{fig:transfer-only-states}
\end{figure}

\subsection{Transfer followed by death}
\label{sec:transfer-death}

Now let each interior particle die at rate \(\mu>0\).  The jumps are
\begin{align}
  (i,j)&\longrightarrow(i+1,j-1) &&\text{at rate }\alpha j,
  \notag\\
  (i,j)&\longrightarrow(i-1,j)   &&\text{at rate }\mu i.
  \label{eq:transfer-death-rates}
\end{align}
The forward equations become
\begin{equation}
  \frac{\dd p_{i,j}}{\dd t}
  =\alpha(j+1)p_{i-1,j+1}+\mu(i+1)p_{i+1,j}
   -(\alpha j+\mu i)p_{i,j}.
  \label{eq:transfer-death-forward}
\end{equation}
Using the two jump increments in the generator gives
\begin{equation}
  \frac{\partial G}{\partial t}
  =\mu(1-x)\frac{\partial G}{\partial x}
   +\alpha(x-y)\frac{\partial G}{\partial y},
  \qquad G(x,y,0)=y^N.
  \label{eq:transfer-death-pde}
\end{equation}

For one initial particle there are three states: exterior \((0,1)\), interior
\((1,0)\), and dead \((0,0)\).  The first two probabilities solve
\begin{equation}
  p_{0,1}'=-\alpha p_{0,1},
  \qquad
  p_{1,0}'=\alpha p_{0,1}-\mu p_{1,0}.
  \label{eq:transfer-death-single-odes}
\end{equation}
When \(\alpha\neq\mu\),
\begin{align}
  p_{0,1}(t)&=\e^{-\alpha t},
  \notag\\
  p_{1,0}(t)&=\frac{\alpha}{\mu-\alpha}
              (\e^{-\alpha t}-\e^{-\mu t}),
  \notag\\
  p_{0,0}(t)&=1-p_{0,1}(t)-p_{1,0}(t).
  \label{eq:transfer-death-states}
\end{align}
At the removable degeneracy \(\alpha=\mu\), the middle expression is replaced by
\begin{equation}
  p_{1,0}(t)=\alpha t\e^{-\alpha t}.
  \label{eq:transfer-death-equal-rates}
\end{equation}
Independence of the initial particles again gives a power:
\begin{equation}
  G_N(x,y,t)
  =\{p_{0,0}(t)+p_{1,0}(t)x+p_{0,1}(t)y\}^{N}.
  \label{eq:G-transfer-death-compact}
\end{equation}
For \(\alpha\neq\mu\), this is
\begin{equation}
  \boxed{\begin{aligned}
  G_N(x,y,t)=\Bigg\{&
  1-\frac{\mu}{\mu-\alpha}\e^{-\alpha t}
    +\frac{\alpha}{\mu-\alpha}\e^{-\mu t} \\
  &+\frac{\alpha}{\mu-\alpha}
      (\e^{-\alpha t}-\e^{-\mu t})x
    +\e^{-\alpha t}y
  \Bigg\}^{N}.
  \end{aligned}}
  \label{eq:G-transfer-death}
\end{equation}

\begin{figure}[htbp]
  \centering
  \includegraphics[width=0.58\textwidth]{abs2.pdf}
  \caption{State probabilities for one particle under transfer followed by death.
  The interior probability rises while transfer dominates and then decays as the
  absorbing dead state accumulates probability.}
  \label{fig:transfer-death-states}
\end{figure}

\subsection{Characteristics tested on the transfer--death PDE}
\label{sec:characteristics-test}

The direct calculation has already supplied \cref{eq:G-transfer-death}.  This makes
\cref{eq:transfer-death-pde} a useful test problem for the method of characteristics:
the method can be developed without leaving its answer unchecked.

Introduce parameters \((s,q,\tau)\) and write
\begin{equation}
  (x,y,t)=(x(s,q,\tau),y(s,q,\tau),t(s,q,\tau)),
\end{equation}
with initial surface
\begin{equation}
  x(s,q,0)=s,
  \qquad y(s,q,0)=q,
  \qquad t(s,q,0)=0.
  \label{eq:characteristic-initial-surface}
\end{equation}
Along a curve with \(s\) and \(q\) fixed, the chain rule gives
\begin{equation}
  \frac{\dd G}{\dd\tau}
  =G_t\frac{\dd t}{\dd\tau}
   +G_x\frac{\dd x}{\dd\tau}
   +G_y\frac{\dd y}{\dd\tau}.
  \label{eq:characteristic-chain-rule}
\end{equation}
Writing \cref{eq:transfer-death-pde} as
\(-G_t+\mu(1-x)G_x+\alpha(x-y)G_y=0\) and matching coefficients gives
\begin{align}
  \frac{\dd G}{\dd\tau}&=0,
  &\frac{\dd t}{\dd\tau}&=-1,
  &\frac{\dd x}{\dd\tau}&=\mu(1-x),
  &\frac{\dd y}{\dd\tau}&=\alpha(x-y).
  \label{eq:transfer-death-characteristics}
\end{align}

The first two equations and \cref{eq:characteristic-initial-surface} give
\begin{equation}
  G(s,q,\tau)=q^N,
  \qquad t=-\tau.
  \label{eq:G-t-characteristics-test}
\end{equation}
Solving the \(x\)-equation gives
\begin{equation}
  1-x=(1-s)\e^{-\mu\tau}.
  \label{eq:x-characteristic-test}
\end{equation}
Substitution into the final equation and use of an integrating factor yield, for
\(\alpha\neq\mu\),
\begin{equation}
  y=1-\frac{\alpha}{\alpha-\mu}(1-s)\e^{-\mu\tau}
   +\left\{q+\frac{\alpha}{\alpha-\mu}(1-s)-1\right\}
    \e^{-\alpha\tau}.
  \label{eq:y-characteristic-test}
\end{equation}
The equal-rate case follows by taking the limit \(\mu\to\alpha\).

It remains to return from \((s,q,\tau)\) to \((x,y,t)\).  Rearranging
\cref{eq:y-characteristic-test} for \(q\), then substituting
\(1-s=(1-x)\e^{\mu\tau}\) from \cref{eq:x-characteristic-test} and
\(\tau=-t\), gives
\begin{equation}
\begin{aligned}
  q={}&1-\frac{\mu}{\mu-\alpha}\e^{-\alpha t}
    +\frac{\alpha}{\mu-\alpha}\e^{-\mu t}
    +\frac{\alpha}{\mu-\alpha}
       (\e^{-\alpha t}-\e^{-\mu t})x
    +\e^{-\alpha t}y.
\end{aligned}
  \label{eq:q-recovered-test}
\end{equation}
Finally \(G=q^N\) by \cref{eq:G-t-characteristics-test}, reproducing
\cref{eq:G-transfer-death}.  The characteristic orientation, signs and initial
surface have therefore all been checked against an independent solution.

\subsection{Transfer, birth and death}
\label{sec:transfer-birth-death}

Add linear birth of interior particles at per-capita rate \(\beta>0\).  The jump
rates are now
\begin{align}
  (i,j)&\longrightarrow(i+1,j-1) &&\text{at rate }\alpha j,
  \notag\\
  (i,j)&\longrightarrow(i+1,j)   &&\text{at rate }\beta i,
  \notag\\
  (i,j)&\longrightarrow(i-1,j)   &&\text{at rate }\mu i.
  \label{eq:abd-rates}
\end{align}
The corresponding forward equations are
\begin{multline}
  \frac{\dd p_{i,j}}{\dd t}
  =\alpha(j+1)p_{i-1,j+1}
   +\mu(i+1)p_{i+1,j}
   +\beta(i-1)p_{i-1,j}\\
   -(\alpha j+\mu i+\beta i)p_{i,j}.
  \label{eq:abd-forward}
\end{multline}
The generator applied to \(x^iy^j\) gives the first-order PDE
\begin{equation}
  \frac{\partial G}{\partial t}
  =\{\mu-(\mu+\beta)x+\beta x^2\}\frac{\partial G}{\partial x}
   +\alpha(x-y)\frac{\partial G}{\partial y},
  \qquad G(x,y,0)=y^N.
  \label{eq:abd-pde}
\end{equation}
Birth makes infinitely many interior states reachable from one transferred
particle.  Independence still implies \(G_N=G_1^N\), but \(G_1\) is now an infinite
series and must itself be determined.  We solve \cref{eq:abd-pde} directly.

\subsubsection{The characteristic system}

Use the same initial surface as in \cref{eq:characteristic-initial-surface}.  The
characteristic equations are
\begin{align}
  \frac{\dd G}{\dd\tau}&=0,
  &\frac{\dd t}{\dd\tau}&=-1,
  \notag\\
  \frac{\dd x}{\dd\tau}
    &=\mu-(\mu+\beta)x+\beta x^2,
  &\frac{\dd y}{\dd\tau}&=\alpha(x-y).
  \label{eq:abd-characteristics}
\end{align}
Thus \(G=q^N\) and \(t=-\tau\).  The \(x\)-equation factorises without a
quadratic surd:
\begin{equation}
  \mu-(\mu+\beta)x+\beta x^2
  =\beta(x-1)(x-\rho),
  \qquad
  \rho:=\frac{\mu}{\beta}.
  \label{eq:abd-factorisation}
\end{equation}
Put
\begin{equation}
  b:=\beta-\mu=\beta(1-\rho).
  \label{eq:abd-net-growth}
\end{equation}
When \(b\neq0\), separation of variables and the condition \(x(s,q,0)=s\)
give
\begin{equation}
  \frac{x-1}{x-\rho}
  =\frac{s-1}{s-\rho}\e^{b\tau}.
  \label{eq:abd-x-ratio}
\end{equation}
Writing
\begin{equation}
  Z=\frac{s-1}{s-\rho},
  \qquad
  \zeta(\tau)=Z\e^{b\tau},
  \label{eq:abd-zeta}
\end{equation}
we may solve \cref{eq:abd-x-ratio} as
\begin{equation}
  x(\tau)
  =\frac{1-\rho\zeta(\tau)}{1-\zeta(\tau)}
  =1+(1-\rho)\frac{\zeta(\tau)}{1-\zeta(\tau)}.
  \label{eq:abd-x-zeta}
\end{equation}

The last characteristic equation is linear once \(x(\tau)\) is known:
\begin{equation}
  \frac{\dd}{\dd\tau}\{\e^{\alpha\tau}y(\tau)\}
  =\alpha\e^{\alpha\tau}x(\tau).
  \label{eq:abd-y-integrating-factor}
\end{equation}
For \(|\zeta|<1\), expand the last quotient in
\cref{eq:abd-x-zeta}.  Termwise integration gives
\begin{equation}
  \Psi(\zeta)
  :=1+\frac{\alpha}{\beta}
       \sum_{n=1}^{\infty}\frac{\zeta^n}{n+\delta},
  \qquad
  \delta:=\frac{\alpha}{b},
  \label{eq:Psi-series}
\end{equation}
for non-resonant parameter values.  Indeed,
\begin{equation}
  \frac{\dd}{\dd\tau}
  \{\e^{\alpha\tau}\Psi(\zeta(\tau))\}
  =\alpha\e^{\alpha\tau}x(\tau),
  \label{eq:Psi-antiderivative}
\end{equation}
because \(\alpha+nb=b(n+\delta)\) and
\(b/\beta=1-\rho\).  Applying \(y(s,q,0)=q\) to
\cref{eq:abd-y-integrating-factor} therefore yields
\begin{equation}
  q=\e^{\alpha\tau}\{y-\Psi(\zeta(\tau))\}+\Psi(Z).
  \label{eq:abd-q-characteristic}
\end{equation}

To return to the original variables, set
\begin{equation}
  W=\frac{x-1}{x-\rho}.
  \label{eq:abd-W}
\end{equation}
Then \(\zeta(\tau)=W\), \(Z=W\e^{bt}\) and \(\tau=-t\).  Since
\(G=q^N\), we obtain the promised generating function.

\begin{proposition}[Transfer--birth--death generating function]
\label{prop:abd-pgf}
Suppose \(\beta\neq\mu\), and define \(b,\rho,\delta,W\) and \(\Psi\) by
\cref{eq:abd-net-growth,eq:Psi-series,eq:abd-W}.  Wherever the displayed series
converges, and elsewhere by continuation along the characteristic,
\begin{equation}
  \boxed{
  G(x,y,t)=
  \left[
    \e^{-\alpha t}\{y-\Psi(W)\}
    +\Psi(W\e^{bt})
  \right]^N.}
  \label{eq:abd-pgf-solution}
\end{equation}
At resonant parameter values the right-hand side is understood by its finite
limit, described below.
\end{proposition}

\begin{proof}
Equations \cref{eq:abd-characteristics,eq:abd-q-characteristic} show that the
quantity inside square brackets is the initial coordinate \(q\) of the
characteristic through \((x,y,t)\).  The initial condition is \(G=q^N\), which
proves \cref{eq:abd-pgf-solution}.  The derivation also proves that it satisfies
\cref{eq:abd-pde}; no uniqueness beyond the usual local uniqueness of the
characteristic initial-value problem is being assumed.
\end{proof}

\subsubsection{A parameter-regular representation}
\label{sec:abd-regular}

The special-function form is compact, but the underlying probability calculation
provides a representation that remains transparent at \(\beta=\mu\), at resonance
and across hypergeometric branch cuts.  Let \(F(x,t)\) be the \pgf{} at time \(t\)
of an interior linear birth--death process started from one particle.  Its backward
equation is
\begin{equation}
  \frac{\partial F}{\partial t}
  =\beta F^2-(\beta+\mu)F+\mu,
  \qquad F(x,0)=x.
  \label{eq:interior-bd-backward}
\end{equation}
For \(\beta\neq\mu\),
\begin{equation}
  \frac{F(x,t)-1}{F(x,t)-\rho}
  =\frac{x-1}{x-\rho}\e^{bt},
  \label{eq:interior-bd-ratio}
\end{equation}
while at \(\beta=\mu\),
\begin{equation}
  F(x,t)
  =1-\frac{1-x}{1+\beta t(1-x)}.
  \label{eq:interior-bd-critical}
\end{equation}

For one initial exterior particle, either transfer has not occurred by time \(t\),
or it occurs at time \(u\in(0,t)\) and leaves an interior birth--death process of
duration \(t-u\).  Conditioning on that transfer time gives
\begin{equation}
  g(x,y,t)
  =\e^{-\alpha t}y
   +\alpha\e^{-\alpha t}
      \int_0^t \e^{\alpha v}F(x,v)\,\dd v.
  \label{eq:abd-convolution}
\end{equation}
Consequently,
\begin{equation}
  G(x,y,t)=g(x,y,t)^N.
  \label{eq:abd-convolution-power}
\end{equation}
This formula is real and finite for \(x,y\in[0,1]\) and all positive rates.  For
\(\beta\neq\mu\), substituting \cref{eq:interior-bd-ratio} into the integral and
using \cref{eq:Psi-antiderivative} recovers
\cref{eq:abd-pgf-solution}.  Thus \cref{eq:abd-convolution} also specifies
unambiguously which continuation of the special-function expression is intended.

Several checks now require almost no algebra.  At \(t=0\),
\cref{eq:abd-convolution} gives \(g=y\), and hence \(G=y^N\).  At
\(x=y=1\), \(F(1,v)=1\), so \(g=1\) and total probability is conserved.
Setting \(x=1\) gives
\begin{equation}
  G(1,y,t)
  =\{\e^{-\alpha t}y+1-\e^{-\alpha t}\}^{N},
  \label{eq:abd-exterior-binomial}
\end{equation}
which is the generating function of
\(Y_t\sim\operatorname{Binomial}(N,\e^{-\alpha t})\).  The exterior marginal is
therefore unaffected by the interior birth and death rates, as the transition
diagram requires.

\subsubsection{Hypergeometric form, domains and resonance}
\label{sec:abd-hypergeometric}

The series in \cref{eq:Psi-series} is a Gauss hypergeometric function.  With
\begin{equation}
  {}_2F_1(a,b;c;z)
  =\sum_{n=0}^{\infty}
    \frac{(a)_n(b)_n}{(c)_n}\frac{z^n}{n!},
  \label{eq:hypergeometric-definition}
\end{equation}
the Pochhammer cancellation
\begin{equation}
  \frac{(\delta)_n}{(\delta+1)_n}
  =\frac{\delta}{\delta+n}
  \label{eq:pochhammer-ratio}
\end{equation}
gives
\begin{equation}
  {}_2F_1(1,\delta;\delta+1;z)
  =\sum_{n=0}^{\infty}\frac{\delta}{\delta+n}z^n.
  \label{eq:hypergeometric-simple-series}
\end{equation}
It follows that
\begin{equation}
  \Psi(z)
  =1+(1-\rho)
   \{{}_2F_1(1,\delta;\delta+1;z)-1\}.
  \label{eq:Psi-hypergeometric}
\end{equation}
A second derivation, based on a general exponential--rational integral, is given in
\cref{app:hypergeometric-integral}.

The elementary series \cref{eq:Psi-series} converges for \(|z|<1\).  If
\(\mu>\beta\), then \(b<0\), \(\rho>1\), and, for \(x\in[0,1]\),
\begin{equation}
  0\leq W\leq\frac{\beta}{\mu}<1,
  \qquad
  0\leq W\e^{bt}\leq W.
\end{equation}
Except at the resonances described next, the series therefore covers the whole
\pgf{} interval.  If \(\beta>\mu\), the arguments may leave the unit disc or cross
the standard hypergeometric cut.  One may then use analytic continuation, tracking
the branch along the characteristic, but the real integral
\cref{eq:abd-convolution} is usually safer.  The apparent singularity of \(W\) at
\(x=\rho\) is also only a coordinate singularity; \(F(\rho,t)=\rho\) and the
continuous limit of \cref{eq:abd-convolution} is regular.

When
\begin{equation}
  \alpha=n(\mu-\beta),
  \qquad n=1,2,\ldots,
  \label{eq:abd-resonance}
\end{equation}
we have \(\delta=-n\), and one denominator in
\cref{eq:Psi-series} vanishes.  The separate \(\Psi\)-terms in
\cref{eq:abd-pgf-solution} then have poles, but their combination does not.  In the
resonant term,
\begin{equation}
  \frac{(W\e^{bt})^n-\e^{-\alpha t}W^n}{n+\delta}
  \label{eq:resonant-cancellation}
\end{equation}
has a finite parameter limit and produces the familiar factor
\(t\e^{-\alpha t}\).  Equivalently, \cref{eq:abd-convolution} contains no singular
denominator and supplies the limit directly.  Resonance is therefore a defect of
the separated antiderivative, not of the generating function.

\subsection{Recovering state probabilities}
\label{sec:coefficient-recovery}

A generating function is useful only insofar as its coefficients can be recovered.
In every model above,
\begin{equation}
  p_{i,j}(t)=[x^iy^j]G(x,y,t)
  =\frac{1}{i!\,j!}
   \left.
   \frac{\partial^{i+j}G}
        {\partial x^i\partial y^j}
   \right|_{x=y=0}.
  \label{eq:bivariate-coefficient}
\end{equation}
For transfer--birth--death, write
\begin{equation}
  h(x,t)
  :=\alpha\e^{-\alpha t}
     \int_0^t\e^{\alpha v}F(x,v)\,\dd v,
  \qquad
  g(x,y,t)=\e^{-\alpha t}y+h(x,t).
  \label{eq:abd-h-definition}
\end{equation}
The finite exterior expansion then separates the two coefficient extractions:
\begin{equation}
  p_{i,j}(t)
  =\binom Nj\e^{-\alpha tj}
   [x^i]h(x,t)^{N-j},
  \qquad 0\leq j\leq N.
  \label{eq:abd-coefficient-formula}
\end{equation}
The remaining univariate coefficient can be obtained by differentiation, by
convolution of truncated series, or by Cauchy's integral formula.  The full
coefficient calculations for all three models are collected in
\cref{app:coefficient-extraction}; the accompanying script implements the
convolution route and checks normalisation.

\FloatBarrier
